% ============================================================================
% Related Work  --  MINIMAL VERSION  (STANDALONE)
%
% COMPILE:  pdflatex related_work_minimal.tex     <-- run it TWICE
%           No BibTeX, no wrapper, no .bib file needed, but you must run
%           pdflatex TWO times: pass 1 writes the reference labels into the
%           .aux file, pass 2 resolves the \cite commands. If you run it only
%           once every citation shows as [?].
%           In TeXShop: hit Typeset twice.
%
% This is the AUTHOR'S OWN draft text and structure, preserved. Changes made:
%   1. Citation placeholders ('mention the paper about ...') filled in.
%   2. ONE addition, agreed 2026-08-02: the paragraph on transformers in ECG
%      being classification rather than forecasting  (marked  % >>> ADDED).
%   3. Light grammatical cleanup only. No restructuring, no new sections.
%
% The fuller proposed restructure is in related_work_full.tex, awaiting review.
%
% Factual issues flagged but DELIBERATELY NOT changed, pending author decision;
% each is marked  % >>> CHECK.
%
% TO MERGE INTO YOUR PAPER LATER: delete everything above \section{Related Work}
% and everything from \begin{thebibliography} onward, then \input what remains.
% ============================================================================

\documentclass[11pt]{article}
\usepackage[T1]{fontenc}
\usepackage[utf8]{inputenc}
\usepackage[margin=1in]{geometry}
\usepackage{amsmath}

\title{Related Work draft}
\author{}
\date{\today}

\begin{document}
\maketitle

% ============================ SECTION STARTS HERE ============================

\section{Related Work}
\label{sec:related}

Chaos forecasting is a challenging field with multiple branches, each with its
own models and data analysis. Consider the electrocardiogram (ECG), which has
progressed a great deal since the early work on nonlinear cardiac
prediction~\cite{babloyantz1988}, through the breakthrough brought by recurrent
architectures such as the LSTM~\cite{yildirim2018}, to the deep convolutional
and residual models that define current practice~\cite{hannun2019,ribeiro2020}.
% >>> CHECK: 'babloyantz1988' is the earliest defensible nonlinear-ECG anchor,
% >>> not a forecasting paper as such. If you want a true "first ECG
% >>> forecasting paper", it needs to be identified; the lineage is murky.
In parallel came the emergence of reservoir computing~\cite{jaeger2004}, which
is used to enhance prediction at a much lower energy cost: next-generation
variants train in milliseconds from only a few hundred
samples~\cite{gauthier2021}, and in cardiac work specifically, echo state
networks and their physics-informed hybrids forecast $15$--$20$ action
potentials with action-potential-duration error below
$5$~ms~\cite{shahi2021,shahi2022}. Most recently, transformer-based models have
been applied to the ECG~\cite{xia2023,zhou2024}.

% >>> ADDED (agreed 2026-08-02) ==============================================
We investigated how transformer models have been used for the ECG. The work we
found applies them to classification, principally hybrid
convolutional--transformer classifiers for inter-patient arrhythmia
recognition~\cite{xia2023,zhou2024}. So far, however, we have found no work that
uses a transformer to forecast the future waveform, which leaves the field open
to a novel contribution.
% >>> END ADDED ==============================================================

According to the literature the results are still limited, particularly over
long horizons. One reason lies in the nature of the data itself, which is
chaotic, and which therefore limits forecasting intrinsically.
% >>> CHECK: whether the ECG is *low-dimensional deterministic chaos* is
% >>> disputed (see Paper_to_be_published/README.md section 7). Your own
% >>> surrogate tests fail on 3 of 6 records. Suggested safer wording:
% >>> "...lies in the nature of the data itself, which is nonlinear,
% >>> non-stationary and multiscale, and whose chaoticity remains debated."
% >>> Left as you wrote it, pending your decision.

\medskip

We therefore turned to another domain, weather, and investigated it. Weather
prediction rests on physical models, namely numerical weather prediction
systems, of which general circulation models (GCMs) are the climate-scale
counterpart. These models
% >>> CHECK: operational *weather* forecasting is NWP (e.g. ECMWF IFS, in HRES
% >>> and ENS configurations). GCM is the climate-modelling term. NeuralGCM is
% >>> named for the GCM lineage because it also runs climate simulations.
which integrate the governing equations directly but require enormous
computational power~\cite{kochkov2024}. This motivated the field to ask whether
machine learning could deliver comparable predictions at some trade-off, whether
in exact pointwise values or in long-term behaviour. That is what we see in
NeuralGCM~\cite{kochkov2024}, which forecasts for roughly one week at a quality
comparable to the physical model, but far faster and at far lower energy cost.
% >>> CHECK: NeuralGCM is evaluated to **15 days** for weather, and separately
% >>> runs **40-year** climate simulations (22 of 37 initialisations stable for
% >>> the full period). It also *beats* ECMWF-ENS on ensemble-mean RMSE and
% >>> CRPS across nearly all variables and leads, rather than merely matching
% >>> it. Throughput: 70,000 simulated days/day on one accelerator vs 19 days
% >>> on 13,824 CPU cores. Left as you wrote it, pending your decision.
Three-dimensional transformer models have also been used~\cite{bi2023}, but so
far the prediction horizon remains on the order of one week.
% >>> CHECK: true for Pangu-Weather specifically (evaluated to 7 days), but
% >>> GraphCast reaches 10 days, ClimaX 14, GenCast 15. Either name Pangu
% >>> explicitly or adjust the claim.

\medskip

\noindent\emph{[Remaining sections to be written.]}

% ============================= SECTION ENDS HERE =============================

% Embedded bibliography: no BibTeX required. If you later switch to BibTeX,
% delete this block and use \bibliography{related_work} instead; the same
% keys are defined in related_work.bib.
\begin{thebibliography}{99}

\bibitem{babloyantz1988}
A.~Babloyantz and A.~Destexhe.
\newblock Is the normal heart a periodic oscillator?
\newblock {\em Biological Cybernetics}, 58(3):203--211, 1988.

\bibitem{yildirim2018}
{\"O}.~Y{\i}ld{\i}r{\i}m.
\newblock A novel wavelet sequence based on deep bidirectional {LSTM} network
  model for {ECG} signal classification.
\newblock {\em Computers in Biology and Medicine}, 96:189--202, 2018.

\bibitem{hannun2019}
A.~Y. Hannun, P.~Rajpurkar, M.~Haghpanahi, et~al.
\newblock Cardiologist-level arrhythmia detection and classification in
  ambulatory electrocardiograms using a deep neural network.
\newblock {\em Nature Medicine}, 25:65--69, 2019.

\bibitem{ribeiro2020}
A.~H. Ribeiro, M.~H. Ribeiro, G.~M.~M. Paix{\~a}o, et~al.
\newblock Automatic diagnosis of the 12-lead {ECG} using a deep neural network.
\newblock {\em Nature Communications}, 11:1760, 2020.

\bibitem{jaeger2004}
H.~Jaeger and H.~Haas.
\newblock Harnessing nonlinearity: Predicting chaotic systems and saving energy
  in wireless communication.
\newblock {\em Science}, 304(5667):78--80, 2004.

\bibitem{gauthier2021}
D.~J. Gauthier, E.~Bollt, A.~Griffith, and W.~A.~S. Barbosa.
\newblock Next generation reservoir computing.
\newblock {\em Nature Communications}, 12:5564, 2021.

\bibitem{shahi2021}
S.~Shahi, C.~D. Marcotte, C.~J. Herndon, F.~H. Fenton, Y.~Shiferaw, and
  E.~M. Cherry.
\newblock Long-time prediction of arrhythmic cardiac action potentials using
  recurrent neural networks and reservoir computing.
\newblock {\em Frontiers in Physiology}, 12:734178, 2021.

\bibitem{shahi2022}
S.~Shahi, F.~H. Fenton, and E.~M. Cherry.
\newblock Prediction of chaotic time series using recurrent neural networks and
  reservoir computing techniques: A comparative study.
\newblock {\em Machine Learning with Applications}, 8:100300, 2022.

\bibitem{xia2023}
Y.~Xia, Y.~Xiong, and K.~Wang.
\newblock A transformer model blended with {CNN} and denoising autoencoder for
  inter-patient {ECG} arrhythmia classification.
\newblock {\em Biomedical Signal Processing and Control}, 86:105271, 2023.

\bibitem{zhou2024}
Y.~Zhou and X.~Wang.
\newblock Heartbeat classification by combining multi-branch convolutional
  neural networks and transformer.
\newblock {\em iScience}, 27(3):109307, 2024.

% [removed at author's request; restore if needed]
% \bibitem{busia2024}
% P.~Busia, M.~A. Scrugli, V.~J.~B. Jung, L.~Benini, and P.~Meloni.
% \newblock A tiny transformer for low-power arrhythmia classification on
%   microcontrollers.
% \newblock {\em IEEE Transactions on Biomedical Circuits and Systems}, 2024.

% [removed at author's request; restore if needed]
% \bibitem{mckeen2025ecgfm}
% K.~McKeen et~al.
% \newblock {ECG-FM}: An open electrocardiogram foundation model.
% \newblock {\em JAMIA Open}, 8(5), 2025.
% \newblock [author list unverified; check before submission].

% [uncited after the transformer paragraph was shortened; restore if needed]
% \bibitem{fade2025}
% N.~Rold{\'a}n et~al.
% \newblock {FADE}: Forecasting for anomaly detection on {ECG}.
% \newblock {\em Computer Methods and Programs in Biomedicine}, 267:108780, 2025.
% \newblock [author list unverified; check before submission].

% [uncited after the transformer paragraph was shortened; restore if needed]
% \bibitem{zhang2025parroting}
% Y.~Zhang and W.~Gilpin.
% \newblock Context parroting: A simple but tough-to-beat baseline for foundation
%   models in scientific machine learning.
% \newblock arXiv:2505.11349, 2025. ICLR 2026.

\bibitem{kochkov2024}
D.~Kochkov, J.~Yuval, I.~Langmore, et~al.
\newblock Neural general circulation models for weather and climate.
\newblock {\em Nature}, 632(8027):1060--1066, 2024.

\bibitem{bi2023}
K.~Bi, L.~Xie, H.~Zhang, X.~Chen, X.~Gu, and Q.~Tian.
\newblock Accurate medium-range global weather forecasting with {3D} neural
  networks.
\newblock {\em Nature}, 619:533--538, 2023.

\end{thebibliography}

\end{document}
